\documentclass[12pt,a4paper]{article}

% -----------------------------
% Packages
% -----------------------------
\usepackage[utf8]{inputenc}
\usepackage[T1]{fontenc}
\usepackage[spanish]{babel}
\usepackage{geometry}
\usepackage{amsmath, amssymb, amsthm}
\usepackage{graphicx}
\usepackage{booktabs}
\usepackage{array}
\usepackage{longtable}
\usepackage{multirow}
\usepackage{hyperref}
\usepackage{xcolor}
\usepackage{enumitem}
\usepackage{float}
\usepackage{caption}
\usepackage{subcaption}
\usepackage{fancyhdr}
\usepackage{setspace}
\usepackage{csquotes}
\usepackage{listings}

% -----------------------------
% Page layout
% -----------------------------
\geometry{
    left=2.5cm,
    right=2.5cm,
    top=2.5cm,
    bottom=2.5cm
}
\setstretch{1.15}
\hypersetup{
    colorlinks=true,
    linkcolor=blue,
    citecolor=blue,
    urlcolor=blue,
    pdftitle={Informe Técnico - Proyecto Final},
    pdfauthor={Tu nombre}
}

% -----------------------------
% Header / Footer
% -----------------------------
\pagestyle{fancy}
\fancyhf{}
\lhead{Proyecto Final de Inteligencia Artificial y Simulación}
\rhead{Informe Técnico}
\cfoot{\thepage}

% -----------------------------
% Metadata
% -----------------------------
\title{Informe Técnico del Proyecto Final\\\large Priorización de Casos Clínicos Bajo Recursos Limitados}
\author{Tu nombre}
\date{\today}

\begin{document}

\maketitle
\tableofcontents
\newpage

% =========================================================
% 1. Introducción / Contexto (opcional pero recomendable)
% =========================================================
\section{Introducción}

La gestión eficiente de recursos hospitalarios constituye uno de los principales desafíos en los sistemas de salud, 
especialmente en contextos donde la disponibilidad de camas, medicamentos, equipos médicos y personal especializado es limitada. 
En estas situaciones, la decisión sobre qué pacientes deben ser atendidos primero puede tener un impacto significativo tanto 
en la cantidad de vidas salvadas como en la calidad de la atención brindada.

Este trabajo aborda el problema de la priorización de casos clínicos bajo restricciones de recursos mediante 
la combinación de técnicas de Inteligencia Artificial y Simulación. 
El objetivo consiste en diseñar un sistema capaz de determinar dinámicamente el orden de atención de los pacientes 
que llegan a un hospital, considerando la evolución temporal de su estado de salud, la disponibilidad de recursos y 
la incertidumbre inherente a los procesos clínicos.

% =========================================================
% 2. Descripción del problema
% =========================================================
\section{Descripción del problema}

Los sistemas hospitalarios operan frecuentemente bajo restricciones de recursos que impiden atender de manera inmediata a 
todos los pacientes que requieren asistencia médica. 
La disponibilidad limitada de camas, medicamentos, equipos y personal especializado obliga a establecer 
mecanismos de priorización que permitan utilizar dichos recursos de la forma más eficiente posible.

En este contexto, se considera un hospital que recibe pacientes a lo largo del tiempo. 
Cada paciente presenta un estado de salud inicial, una velocidad de deterioro cuando no recibe tratamiento y 
una capacidad de recuperación cuando es admitido y atendido. 
Además, cada caso requiere determinados recursos médicos cuya disponibilidad es limitada dentro del hospital.

La admisión de pacientes está restringida tanto por la capacidad física del hospital como por la disponibilidad de 
los recursos necesarios para llevar a cabo el tratamiento. 
Cuando no es posible admitir inmediatamente a todos los pacientes que han llegado, 
estos permanecen en una cola de espera donde continúan evolucionando clínicamente. 
Durante este período, su estado de salud puede deteriorarse progresivamente, incrementando el riesgo de complicaciones o 
fallecimiento.

El problema consiste en determinar dinámicamente el orden en que los pacientes deben ser admitidos para recibir tratamiento. 
Dicha decisión debe considerar simultáneamente la evolución temporal de los pacientes, 
la disponibilidad cambiante de recursos y la incertidumbre asociada a la respuesta de cada individuo al tratamiento.

La calidad de una política de priorización no puede evaluarse únicamente a partir de la información inicial de los pacientes, 
ya que el resultado final depende de la interacción entre múltiples factores que evolucionan a lo largo del tiempo. 
Por esta razón, la evaluación de una solución requiere simular el comportamiento completo del sistema hospitalario, 
incluyendo la llegada de pacientes, la utilización de recursos, la evolución clínica de los casos y 
la liberación de capacidad hospitalaria.

El objetivo del sistema es encontrar estrategias de admisión que maximicen el impacto total de la atención médica. 
Para ello se adopta un criterio de optimización lexicográfico compuesto por dos objetivos jerárquicos. 
En primer lugar, se busca maximizar la cantidad de pacientes que sobreviven al finalizar la simulación. 
En segundo lugar, entre todas las soluciones que producen la misma cantidad de supervivientes, 
se selecciona aquella que minimiza el sufrimiento total experimentado por los pacientes.

El sufrimiento individual de un paciente se define como el menor nivel de salud alcanzado durante toda su evolución 
dentro del sistema. 
A partir de esta definición, el sufrimiento total se obtiene agregando el sufrimiento individual de todos 
los pacientes considerados en la simulación. 
De esta forma, el sistema no solo busca salvar la mayor cantidad posible de vidas, 
sino también evitar que los pacientes alcancen estados críticos de deterioro mientras esperan o reciben atención.

% =========================================================
% 3. Modelado formal
% =========================================================
\section{Modelado formal}

\subsection{Conjunto de pacientes}

Sea

\[
P = \{p_1,p_2,\ldots,p_n\}
\]

el conjunto de pacientes que pueden llegar al hospital durante la simulación.

Cada paciente \(p_i\) se caracteriza mediante la tupla:

\[
p_i =
\left(
h_i^0,
d_i,
r_i,
R_i
\right)
\]

donde:

\begin{itemize}
    \item \(h_i^0 \in [1,10]\) es el nivel de salud inicial.
    \item \(d_i \in [1,10]\) es la tasa de deterioro.
    \item \(r_i \in [1,10]\) es la tasa de recuperación.
    \item \(R_i\) es el conjunto de recursos requeridos para su tratamiento.
\end{itemize}

Durante la simulación cada paciente puede encontrarse en uno de los siguientes estados:

\[
\{\texttt{pending},
\texttt{waiting},
\texttt{admitted},
\texttt{discharged},
\texttt{deceased}\}
\]

correspondientes a espera, ingreso, alta médica y fallecimiento.

\subsection{Recursos del hospital}

Sea

\[
\mathcal{R} =
\{r_1,r_2,\ldots,r_m\}
\]

el conjunto de tipos de recursos disponibles en el hospital.

Cada recurso \(r_j\) posee una disponibilidad inicial:

\[
A_j \ge 0
\]

que representa la cantidad total disponible para ser utilizada durante la simulación.

Además, el hospital dispone de una capacidad máxima de atención:

\[
C \in \mathbb{N}
\]

que indica la cantidad máxima de pacientes que pueden permanecer ingresados simultáneamente.

La capacidad y disponibilidad de recursos se generan sintéticamente a partir de las necesidades agregadas de los pacientes utilizando factores de holgura aleatorios que permiten construir escenarios con distintos niveles de escasez.

\subsection{Variables de estado}

Sea \(t\) un instante discreto de simulación.

El estado de salud de un paciente \(p_i\) en el instante \(t\) se representa mediante:

\[
h_i(t)\in[1,10]
\]

donde:

\[
1
\]

representa el peor estado posible y corresponde al fallecimiento, mientras que

\[
10
\]

representa la recuperación completa y genera el alta médica.

Para cada paciente se registra además:

\[
m_i =
\min_t h_i(t)
\]

que representa el nivel mínimo de salud alcanzado durante toda su evolución.

Esta magnitud será utilizada posteriormente para medir el sufrimiento individual.

\subsection{Evolución temporal de los pacientes}

La simulación avanza mediante pasos discretos de tiempo.

Los pacientes llegan al sistema de acuerdo con un proceso de llegadas exponencial. Si \(T_i\) representa el instante de llegada del paciente \(p_i\), entonces los tiempos entre llegadas siguen una distribución exponencial:

\[
\Delta T_i
\sim
Exp(\lambda)
\]

donde \(\lambda\) está determinado por el tiempo medio entre llegadas configurado en la simulación.

\subsubsection{Paciente en espera}

Cuando un paciente no ha sido admitido en el hospital, su estado de salud disminuye según:

\[
h_i(t+1)
=
h_i(t)
-
D_i(t)
\]

donde

\[
D_i(t)
=
\frac{
\min(10,\max(1,d_i \varepsilon_t))
}{10}
\]

y

\[
\varepsilon_t
\sim
\mathcal{N}(1,\sigma^2)
\]

es un factor aleatorio que introduce incertidumbre clínica.

Si

\[
h_i(t)\le 1
\]

el paciente fallece.

\subsubsection{Paciente ingresado}

Cuando un paciente se encuentra ingresado y recibe tratamiento, su salud evoluciona según:

\[
h_i(t+1)
=
h_i(t)
+
I_i(t)
\]

donde

\[
I_i(t)
=
\frac{
\min(10,\max(1,r_i \varepsilon_t))
}{20}
\]

y nuevamente

\[
\varepsilon_t
\sim
\mathcal{N}(1,\sigma^2).
\]

Si

\[
h_i(t)\ge 10
\]

el paciente recibe el alta médica.

\subsection{Función objetivo}

La calidad de una solución se evalúa mediante un criterio lexicográfico compuesto por dos objetivos jerárquicos.

Sea

\[
S
\]

la cantidad total de pacientes que sobreviven al finalizar la simulación.

El objetivo principal consiste en:

\[
\max S
\]

Entre todas las soluciones con igual valor de \(S\), se selecciona aquella que minimiza el sufrimiento total.

El sufrimiento individual del paciente \(p_i\) se define como:

\[
m_i
=
\min_t h_i(t)
\]

y el sufrimiento total como:

\[
F
=
\sum_{i=1}^{n}
(10-m_i)
\]

De esta forma se penalizan las soluciones que permiten que los pacientes alcancen estados críticos de salud.

El problema completo se expresa como:

\[
\max (S,-F)
\]

bajo orden lexicográfico.

\subsection{Restricciones}

Toda solución factible debe satisfacer simultáneamente:

\begin{enumerate}
    \item Restricción de capacidad hospitalaria:

    \[
    |\mathcal{A}(t)| \le C
    \]

    donde \(\mathcal{A}(t)\) es el conjunto de pacientes ingresados en el instante \(t\).

    \item Restricciones de recursos:

    \[
    u_j(t)
    \le
    A_j(t)
    \]

    para todo recurso \(r_j\), donde \(u_j(t)\) representa la cantidad consumida.

    \item Restricción de admisión:

    Un paciente solo puede ser admitido si existe capacidad disponible y todos los recursos requeridos para iniciar su tratamiento se encuentran disponibles.

\end{enumerate}

\subsection{Formulación del problema}

El problema puede interpretarse como la búsqueda de una política de priorización que determine el orden de atención de los pacientes a medida que estos llegan al sistema.

\subsubsection{Definiciones}

Sea

\[
Q(t)
\]

la cola de pacientes en espera en el instante \(t\).

Una política de priorización es una función:

\[
\pi :
Q(t)
\rightarrow
Q'(t)
\]

que reordena la cola antes de cada decisión de admisión.

\subsubsection{Notación}

\begin{center}
\begin{tabular}{ll}
\toprule
Símbolo & Descripción \\
\midrule
\(P\) & Conjunto de pacientes \\
\(C\) & Capacidad hospitalaria \\
\(h_i(t)\) & Salud del paciente \(i\) en el tiempo \(t\) \\
\(d_i\) & Tasa de deterioro \\
\(r_i\) & Tasa de recuperación \\
\(m_i\) & Salud mínima alcanzada \\
\(S\) & Cantidad de supervivientes \\
\(F\) & Sufrimiento total \\
\(Q(t)\) & Cola de espera \\
\(\pi\) & Política de priorización \\
\bottomrule
\end{tabular}
\end{center}

\subsubsection{Interpretación del modelo}

Bajo esta formulación, el problema se convierte en un problema de optimización dinámica y estocástica. Las decisiones de admisión modifican la evolución futura del sistema, mientras que la incertidumbre clínica introduce variabilidad en los resultados obtenidos. Debido a que la calidad de una política depende de la interacción entre llegadas, disponibilidad de recursos y evolución de los pacientes, la evaluación de una solución requiere ejecutar la simulación completa del sistema hospitalario y analizar sus resultados finales.

% =========================================================
% 4. Dataset
% =========================================================
\section{Dataset utilizado}

El dataset utilizado en este trabajo fue construido a partir de una base pública de prescripción de medicamentos 
asociada a enfermedades, y posteriormente transformado en un conjunto de pacientes sintéticos con información estructurada 
y narrativa clínica generada mediante un modelo de lenguaje. 
Este enfoque permitió disponer de casos diversos, realistas y compatibles con el problema de priorización clínica 
bajo recursos limitados.

\subsection{Origen del dataset}

La fuente inicial del proyecto fue descargada desde Kaggle, 
específicamente el dataset \emph{Drug Prescription to Disease Dataset}. 
El archivo original contiene 14\,682 filas y modela relaciones del tipo enfermedad--medicamento, 
en las que cada registro indica qué fármaco ha sido utilizado para tratar una determinada patología.

A partir de ese CSV original se construyó una representación auxiliar denominada \texttt{drug2patient}, 
la cual invierte la relación inicial y agrupa, para cada medicamento, el conjunto de enfermedades asociadas. 
Esta transformación resultó útil para enriquecer la generación de pacientes sintéticos y asegurar coherencia entre 
diagnóstico y tratamiento.

\subsection{Estructura general}

El dataset final no se utiliza directamente en su forma tabular original, 
sino como punto de partida para generar pacientes sintéticos. La estructura general del proceso es la siguiente:

\[
\text{enfermedad + edad + triage} \rightarrow \text{narrativa clínica} \rightarrow \text{extracción JSON} \rightarrow \text{paciente estructurado}
\]

En la implementación, cada paciente se genera primero como un objeto básico con atributos mínimos, 
luego se expande su descripción clínica mediante un LLM y, finalmente, 
se extraen los campos estructurados que alimentan el simulador y el algoritmo de priorización.

\subsection{Campos de los pacientes}

Cada paciente del dataset final contiene los siguientes campos:

\begin{itemize}
    \item \texttt{id}: identificador único del paciente.
    \item \texttt{age}: edad del paciente.
    \item \texttt{diagnosis}: descripción clínica generada automáticamente en lenguaje natural.
    \item \texttt{health\_level}: nivel inicial de salud, en el rango \([1,10]\).
    \item \texttt{deterioration\_rate}: tasa de deterioro clínico, en el rango \([1,10]\).
    \item \texttt{improvement\_rate}: tasa de mejoría esperada con tratamiento, en el rango \([1,10]\).
    \item \texttt{drugs}: lista de medicamentos asociados al caso, con sus dosis numéricas.
\end{itemize}

Además, a nivel interno del sistema, cada paciente queda asociado a los recursos requeridos para su tratamiento, 
lo cual permite modelar la disponibilidad hospitalaria y las restricciones de admisión.

\subsection{Proceso de generación}

El proceso de generación se divide en varias etapas.

En primer lugar, se construye un paciente básico mediante un generador aleatorio controlado. En esta etapa se selecciona:

\begin{itemize}
    \item la edad, usando una distribución beta ajustada al rango \([0,100]\);
    \item el triage, elegido aleatoriamente entre \texttt{Red}, \texttt{Yellow} y \texttt{Green};
    \item la enfermedad principal, tomada desde un conjunto de patologías frecuentes;
    \item los medicamentos asociados a dicha enfermedad, obtenidos mediante una función de mapeo entre enfermedad y tratamiento.
\end{itemize}

Posteriormente, ese paciente básico se envía a un LLM mediante un prompt especializado para 
generar una narrativa clínica realista. 
Esta narrativa incorpora la edad, el triage, la enfermedad y los medicamentos requeridos, 
manteniendo coherencia semántica entre todos los elementos.

A continuación, la narrativa generada se utiliza como entrada de un segundo prompt 
que obliga al LLM a devolver un objeto JSON con los atributos estructurados del paciente. 
En esta etapa el modelo estima el nivel de salud, la tasa de deterioro, la tasa de mejoría y 
la lista de medicamentos con sus dosis.

Finalmente, los datos básicos del paciente, la narrativa clínica y la salida JSON del LLM 
se combinan en un único registro estructurado. Repetido este proceso para múltiples instancias, 
se obtiene el dataset final de pacientes.

\subsection{Justificación de validez}

El dataset construido cumple con las condiciones del problema por varias razones. 
En primer lugar, cada paciente contiene información estructurada y una descripción textual, 
lo que permite que el LLM cumpla un rol funcional real en el sistema y no meramente decorativo. 
En segundo lugar, la presencia de medicamentos y dosis introduce restricciones concretas de recursos, 
generando escenarios donde la priorización no es trivial.

Además, el uso de una base pública de enfermedades y medicamentos reduce la arbitrariedad de los casos sintéticos 
y favorece la coherencia entre diagnóstico y tratamiento. La aleatoriedad controlada en edad, triage y selección 
de patología produce diversidad clínica suficiente para construir instancias difíciles, en las que el algoritmo debe 
decidir entre pacientes con distintos niveles de gravedad, riesgo y demanda de recursos.

Por último, la estructura obtenida es compatible con la simulación de eventos discretos 
y con la función objetivo definida para el proyecto, ya que cada caso puede evaluarse según su evolución temporal 
y el impacto clínico alcanzado bajo limitaciones hospitalarias.

% =========================================================
% 5. Diseño del algoritmo
% =========================================================
\section{Diseño del algoritmo}

\subsection{Visión general}

El algoritmo propuesto se implementa como una simulación de eventos discretos en la que 
el hospital atiende pacientes bajo restricciones de capacidad y recursos. 
La simulación avanza en pasos de tiempo, durante los cuales pueden ocurrir tres procesos principales: 
la llegada de nuevos pacientes, la admisión de pacientes desde la cola de espera y 
la evolución clínica de los pacientes ya ingresados o aún en espera.

La solución se organiza de forma modular para separar claramente el modelo del dominio, 
la dinámica clínica, la simulación temporal y las políticas de priorización. 
Esta separación permite comparar distintas estrategias de atención sin modificar la lógica central del sistema.

\subsection{Arquitectura de la solución}

La implementación se estructura en tres niveles principales.

\begin{itemize}
    \item \textbf{Dominio}: define las entidades fundamentales del problema, principalmente \texttt{Patient}, \texttt{Hospital} y \texttt{Resource}.
    \item \textbf{Dinámica}: modela la evolución de la salud y el consumo progresivo de recursos mediante funciones estocásticas.
    \item \textbf{Simulación}: coordina el paso del tiempo, las llegadas, la cola de espera, la admisión y la evolución clínica.
    \item \textbf{Políticas}: implementa los distintos criterios para reordenar pacientes antes de la admisión.
\end{itemize}

La clase abstracta \texttt{BaseSimulator} encapsula la lógica común de simulación, 
mientras que \texttt{BasicSimulator} e \texttt{IntelligentSimulator} concretan dos modos de ejecución distintos. 
El primero asume que todos los pacientes ya están en la cola inicial; el segundo simula llegadas en el tiempo y 
reordena la cola usando una política de priorización.

\subsection{Procesamiento de llegadas}

En el simulador inteligente, los pacientes no poseen un tiempo de llegada almacenado en su estructura de datos. 
En su lugar, el propio simulador genera los instantes de llegada de forma interna usando una distribución exponencial, 
lo que permite representar un proceso estocástico de arribos.

Al inicio de la simulación se construye una lista de llegadas programadas. 
En cada paso temporal, el simulador revisa si existen pacientes cuyo tiempo de llegada ha sido alcanzado y 
los incorpora a la cola de espera. 
De esta forma, la llegada de pacientes se modela como un evento discreto independiente de la definición del 
paciente en el dominio.

Cuando la política es \texttt{None}, no se simulan llegadas; 
todos los pacientes se consideran disponibles desde el comienzo y se insertan directamente en la cola inicial.

\subsection{Construcción y reordenamiento de la cola}

La cola de espera contiene a los pacientes que ya han llegado al hospital pero aún no han sido admitidos. 
En cada iteración de la simulación, esta cola puede reordenarse según la política activa.

La interfaz \texttt{PriorityPolicy} define un método \texttt{order\_patients} que recibe la cola actual, 
el estado del hospital y, opcionalmente, la lista de nuevos pacientes llegados en ese instante. 
La política devuelve una nueva secuencia ordenada que será utilizada para intentar admisiones.

Se implementan varias políticas de comparación, entre ellas:

\begin{itemize}
    \item \textbf{RandomPolicy}: orden aleatorio de la cola.
    \item \textbf{SeverityPolicy}: prioridad al paciente con menor nivel de salud.
    \item \textbf{NegligentPolicy}: prioridad inversa a la severidad.
    \item \textbf{SimAnnealPolicy}: búsqueda por recocido simulado sobre permutaciones.
    \item \textbf{ILSPolicy}: búsqueda local iterada.
    \item \textbf{SAILSPolicy}: combinación híbrida de búsqueda local iterada y recocido simulado.
\end{itemize}

De esta manera, el sistema puede comparar estrategias simples contra heurísticas y metaheurísticas más sofisticadas.

\subsection{Política de admisión}

Una vez ordenada la cola, el hospital intenta admitir siempre al primer paciente. 
Si dicho paciente no puede ser admitido por falta de capacidad o por insuficiencia de recursos, 
entonces no se considera a ningún otro paciente durante ese ciclo.

Esta decisión simplifica el mecanismo de admisión y refuerza el carácter secuencial del problema. 
En particular, la cola no se trata como una lista de candidatos intercambiables en cada instante, 
sino como una secuencia priorizada cuya cabeza representa la mejor opción disponible.

La admisión en sí no descuenta de inmediato el consumo progresivo de recursos. 
Solo verifica que el paciente puede ser incorporado al hospital. 
El consumo real de recursos ocurre posteriormente, durante la evolución clínica.

\subsection{Evolución clínica y consumo de recursos}

Cada paciente ingresado evoluciona durante la simulación de forma estocástica. 
La salud del paciente puede mejorar si recibe tratamiento o deteriorarse si no dispone de los recursos necesarios.

Las funciones \texttt{compute\_deterioration} y \texttt{compute\_improvement} introducen ruido aleatorio 
mediante una distribución normal centrada en 1.0. A partir de estas funciones se actualiza el estado del 
paciente de la siguiente manera:

\begin{itemize}
    \item si el paciente está en espera, su salud disminuye;
    \item si el paciente está ingresado y recibe tratamiento, su salud mejora;
    \item si el paciente alcanza el nivel mínimo de salud, fallece;
    \item si alcanza el nivel máximo de salud, recibe el alta.
\end{itemize}

El consumo de medicamentos es progresivo. Cada paciente posee una cantidad total requerida de recursos, 
y en cada paso temporal se consume solo la parte no administrada previamente. 
Este consumo se calcula a partir de la cantidad pendiente y del estado actual del paciente, 
de modo que el tratamiento se distribuye a lo largo del tiempo en lugar de consumirse de una sola vez.

Si el hospital no dispone del recurso necesario para cubrir el consumo de un paciente en un instante dado, 
el tratamiento no puede aplicarse completamente y el paciente puede deteriorarse. 
Este mecanismo introduce competencia real entre pacientes por recursos limitados.

\subsection{Estrategia de búsqueda u optimización}

El objetivo del algoritmo no es únicamente simular el sistema, 
sino encontrar una política de atención que produzca una solución cercana al óptimo. 
Para ello se emplean diferentes estrategias de priorización.

Las políticas simples actúan como baselines. Ordenan a los pacientes por criterios directos como salud actual o aleatoriedad. 
Las políticas avanzadas utilizan búsquedas sobre el espacio de permutaciones, 
donde cada solución candidata representa un posible orden de atención.

En el caso de \texttt{SimAnnealPolicy}, se modela una permutación como un estado de un problema de optimización y 
se exploran vecinos generados mediante intercambios aleatorios de posiciones. 
La función de energía se define como el negativo de la evaluación de la secuencia, 
de forma que el recocido simulado busca minimizar dicha energía, es decir, maximizar la calidad de la solución.

En el caso de \texttt{ILSPolicy}, se modela una permutación como la secuencia ordenada de atención de los pacientes 
y se emplea un enfoque heurístico iterativo estructurado en tres fases. 
En primer lugar, se utiliza un mecanismo de inserción codiciosa (\texttt{greedy insertion}) para posicionar 
de manera inicial a cada nuevo paciente conforme arriba al sistema. Posteriormente, se aplica una búsqueda 
local basada en el algoritmo de escalada de colinas (\texttt{hill climbing}) mediante el intercambio de posiciones 
adyacentes para refinar el ordenamiento hasta converger en un óptimo local. Para escapar de estos máximos locales y 
explorar nuevas regiones del espacio de soluciones, el algoritmo introduce una fase de perturbación aleatoria —o sacudida— 
en la secuencia, repitiendo el ciclo de optimización local de forma controlada. 
La función de evaluación determina la calidad de cada ordenamiento, orientando al \texttt{ILS} hacia la configuración que maximice 
el rendimiento global del sistema de salud y, por consiguiente, la cantidad de pacientes salvados.

La política \texttt{SAILSPolicy} combina una mejora local iterada con recocido simulado, 
lo que permite refinar soluciones iniciales y escapar de óptimos locales. 
Esta estrategia híbrida resulta especialmente útil cuando el espacio de búsqueda es grande y la evaluación 
exacta de todas las permutaciones es inviable.

\subsection{Complejidad computacional}

La complejidad del algoritmo depende del número de pacientes, del número de pasos de simulación y del tipo de política utilizada.

\subsubsection{Componentes deterministas}

Los componentes deterministas incluyen la admisión del primer paciente factible de la cola, 
la actualización del estado del hospital y la verificación de recursos. 
Si la cola contiene \(n\) pacientes y cada paciente requiere revisar un número \(m\) de recursos, 
una admisión puede realizarse en tiempo proporcional a \(O(m)\), 
mientras que el reordenamiento de la cola depende de la política seleccionada.

\subsubsection{Componentes estocásticos}

Los componentes estocásticos aparecen en la generación de llegadas, en la evolución clínica y en las funciones de deterioro y 
mejoría. Aunque no incrementan necesariamente la complejidad asintótica, sí obligan a repetir las simulaciones varias veces 
para obtener resultados estables y representativos.

\subsubsection{Estrategia de mejora local}

Cuando se utilizan metaheurísticas, la complejidad total aumenta por la necesidad de evaluar múltiples permutaciones candidatas. 
Si cada evaluación requiere una simulación completa, el costo computacional puede ser elevado. 
Para mitigar este problema, la búsqueda local se aplica sobre instancias moderadas o mediante un número controlado de iteraciones.

En términos prácticos, esto permite explorar soluciones cercanas al óptimo sin realizar una búsqueda exhaustiva del 
espacio de permutaciones, que sería factorial en el número de pacientes. La combinación de simulación, heurísticas y 
metaheurísticas ofrece un compromiso adecuado entre calidad de la solución y costo computacional.

% =========================================================
% 6. Rol del LLM
% =========================================================
\section{Rol del LLM en el sistema}

El modelo de lenguaje desempeña un papel fundamental en la construcción del dataset utilizado por el sistema. 
Su función principal consiste en transformar información clínica básica y estructurada en descripciones médicas realistas, 
y posteriormente extraer de dichas descripciones los atributos cuantitativos necesarios para la simulación.

En una primera etapa, el LLM recibe como entrada la edad del paciente, el nivel de triage, la enfermedad principal y 
los medicamentos asociados al tratamiento. 
A partir de esta información genera una narrativa clínica coherente en lenguaje natural que describe el estado del paciente y 
su tratamiento.

Posteriormente, esa narrativa es procesada nuevamente por el LLM para obtener una representación estructurada en formato JSON. 
En esta fase el modelo estima el nivel inicial de salud (\texttt{health\_level}), 
la tasa de deterioro (\texttt{deterioration\_rate}), la tasa de mejoría (\texttt{improvement\_rate}) y 
extrae los medicamentos junto con sus dosis correspondientes.

De esta forma, el LLM actúa como un mecanismo de enriquecimiento semántico y extracción de conocimiento, 
permitiendo convertir información clínica básica en instancias completas compatibles con el simulador y 
con el problema de optimización planteado.

% =========================================================
% 7. Metodología experimental
% =========================================================
\section{Metodología experimental}

El proyecto fue evaluado mediante simulaciones repetidas sobre distintas instancias sintéticas del problema, 
con el objetivo de comparar el comportamiento de varias políticas de priorización bajo diferentes niveles de demanda y 
disponibilidad de recursos. 
La metodología experimental se diseñó para verificar no solo la calidad de la solución obtenida, 
sino también su costo computacional y su estabilidad frente a cambios en la instancia.

\subsection{Diseño experimental}

El diseño experimental se basó en la comparación de varias configuraciones de atención sobre un mismo conjunto de pacientes 
y un mismo hospital sintético. 
Para cada instancia se ejecutó el simulador completo con distintas políticas, 
y luego se compararon los resultados obtenidos según las métricas definidas para el problema.

Las políticas evaluadas fueron:

\begin{itemize}
    \item \textbf{Random Policy}: orden aleatorio de los pacientes.
    \item \textbf{Severity Policy}: priorización por menor nivel de salud.
    \item \textbf{Negligent Policy}: priorización inversa a la severidad.
    \item \textbf{Simulated Annealing Policy}: optimización por recocido simulado.
    \item \textbf{Iterated Local Search Policy}: mejora local iterada.
    \item \textbf{Hybrid Policy}: combinación de recocido simulado e iterated local search.
\end{itemize}

Estas configuraciones permitieron comparar estrategias muy simples contra métodos metaheurísticos más costosos, 
pero potencialmente más cercanos al óptimo.

\subsection{Instancias de prueba}

Se generaron varias instancias de prueba con tamaños distintos, a fin de evaluar el comportamiento del sistema 
tanto en escenarios pequeños como en escenarios más exigentes. En particular, se realizaron experimentos con:

\begin{itemize}
    \item 45 pacientes;
    \item 90 pacientes;
    \item 180 pacientes.
\end{itemize}

Además, para analizar el efecto de la escasez de recursos y la severidad promedio de los casos, 
se utilizaron distintas configuraciones del generador de hospitales, 
variando los parámetros \texttt{mean} y \texttt{variance}. 
Esto permitió construir escenarios con mayor o menor holgura en capacidad y recursos.

\subsection{Configuraciones comparadas}

Para cada instancia se ejecutaron las mismas políticas de priorización, manteniendo constantes los datos de entrada y 
variando únicamente el criterio de ordenación o búsqueda. 
En el caso de las metaheurísticas, también se ajustaron sus parámetros internos, como la temperatura inicial, 
el factor de decaimiento y el número máximo de iteraciones.

Los parámetros utilizados en cada corrida se registraron explícitamente para asegurar la reproducibilidad de los experimentos. 
De esta forma, fue posible comparar tanto la calidad de la solución como el tiempo requerido por cada estrategia.

\subsection{Métricas empleadas}

La evaluación de cada política se realizó mediante las siguientes métricas:

\begin{itemize}
    \item \textbf{Cantidad de vidas salvadas}: número de pacientes que alcanzaron el alta médica.
    \item \textbf{Número de muertes}: cantidad de pacientes que fallecieron durante la simulación.
    \item \textbf{Sufrimiento total}: suma de los niveles mínimos de salud alcanzados por todos los pacientes.
    \item \textbf{Tiempo de simulación}: duración total del proceso simulado en instantes discretos.
    \item \textbf{Tiempo de ejecución}: tiempo real de cómputo requerido por cada política.
\end{itemize}

Dado que el objetivo del problema es lexicográfico, la métrica principal fue la cantidad de vidas salvadas; 
en caso de empate, se utilizó el sufrimiento total como criterio de desempate.

\subsection{Procedimiento de evaluación}

Cada experimento siguió el mismo procedimiento:

\begin{enumerate}
    \item Se generó una instancia del problema con un número fijo de pacientes.
    \item Se construyó el hospital sintético con capacidad y recursos determinados por el generador.
    \item Se ejecutó la simulación con cada una de las políticas definidas.
    \item Se registraron las métricas finales de cada corrida.
    \item Se compararon los resultados y se seleccionó la política ganadora según el criterio lexicográfico.
\end{enumerate}

Para las políticas metaheurísticas, la evaluación de una permutación se realizó mediante simulación completa, 
ya que la evolución de los pacientes es estocástica y no puede estimarse con precisión a partir del orden inicial únicamente.

\subsubsection{Parámetros del simulador}

En los experimentos realizados se utilizaron configuraciones con distintos valores del tiempo medio entre llegadas, 
así como diferentes factores de holgura para capacidad y recursos del hospital. 
En particular, se evaluaron escenarios con \texttt{mean} igual a 1.0, 1.1 y 1.2, y varianza igual a 0.0, 0.5 y 0.7, respectivamente. 
Estos valores permitieron estudiar cómo cambia el desempeño de las políticas cuando el sistema pasa de condiciones más holgadas a condiciones más restrictivas.

\subsubsection{Baselines}

Las políticas \texttt{Random} y \texttt{Negligent} funcionaron como referencias de bajo desempeño, 
mientras que \texttt{Severity} representó una heurística clínica intuitiva. 
Las políticas \texttt{Simulated Annealing}, \texttt{Iterated Local Search} y la política híbrida representaron estrategias 
de optimización más avanzadas.

\subsection{Análisis de los resultados}

Los resultados obtenidos muestran un patrón consistente: las políticas metaheurísticas superan en general a 
las heurísticas simples en términos de vidas salvadas y, en muchos casos, también en sufrimiento total.

En la instancia de 90 pacientes, la política híbrida obtuvo el mejor desempeño, alcanzando 25 vidas salvadas, 
65 muertes y un sufrimiento total de 135.0. 
En la instancia de 45 pacientes, nuevamente la política híbrida obtuvo el mayor número de vidas salvadas, 
con 15 pacientes dados de alta. En la instancia de 180 pacientes con parámetros estándar, 
el recocido simulado obtuvo el mejor resultado en cantidad de vidas salvadas, con 47, superando a las políticas simples y 
a la búsqueda local iterada.

Cuando se aumentó la holgura del hospital mediante los parámetros \texttt{mean=1.1} y \texttt{variance=0.5}, 
la política híbrida volvió a ser la mejor, con 87 vidas salvadas. 
En el escenario más favorable de \texttt{mean=1.2} y \texttt{variance=0.7}, la política híbrida también obtuvo el mejor resultado, 
con 89 vidas salvadas y el menor sufrimiento global entre las configuraciones comparadas.

En términos computacionales, las políticas simples fueron mucho más rápidas, 
pero produjeron soluciones de menor calidad. 
Las metaheurísticas mejoraron el desempeño clínico a costa de un mayor tiempo de ejecución, 
especialmente en el caso de \texttt{Iterated Local Search} y de la política híbrida, 
que requieren múltiples evaluaciones por simulación.

En conjunto, los experimentos confirman que la búsqueda metaheurística permite aproximarse mejor al óptimo 
que las reglas heurísticas simples, especialmente en escenarios con mayor complejidad y escasez de recursos.

% =========================================================
% 8. Resultados y análisis
% =========================================================
\section{Resultados y Análisis}

Los ensayos experimentales revelaron un comportamiento consistente entre las estrategias metaheurísticas examinadas. 
Tanto \texttt{Iterated Local Search} como \texttt{Simulated Annealing} fundamentan su éxito en la exploración 
sistemática de vecindarios sobre el espacio de permutaciones de la cola de pacientes, 
logrando refinar progresivamente las soluciones iniciales. En la totalidad de las instancias críticas evaluadas, 
ambos enfoques superaron de manera categórica a las heurísticas basales, 
optimizando simultáneamente las dos métricas fundamentales del sistema: la maximización de vidas salvadas y 
la minimización del índice de sufrimiento acumulado.

\subsection{Evaluación del Rendimiento Global}

El análisis comparativo de los escenarios de prueba permite identificar una clara gradación en la 
efectividad de las políticas. En el extremo inferior, los enfoques ingenuos o estocásticos 
(\texttt{RandomPolicy} y \texttt{NegligentPolicy}) exhibieron el desempeño más deficiente debido a 
su incapacidad para modelar la urgencia clínica. Si bien la incorporación de criterios médicos en 
\texttt{SeverityPolicy} mitiga parcialmente este impacto, su naturaleza determinista (greedy) 
limita su efectividad ante flujos complejos de pacientes. 
En contraposición, \texttt{Simulated Annealing} e \texttt{Iterated Local Search} 
convergieron de manera recurrente a regiones subóptimas de alta calidad, 
demostrando una eficacia equivalente en la explotación del espacio de búsqueda.

\begin{table}[htbp]
\centering
\caption{Clasificación taxonómica de las políticas según su rendimiento}
\label{tab:taxonomia_politicas}
\begin{tabular}{lll}
\hline
\textbf{Categoría} & \textbf{Políticas} & \textbf{Eficiencia Relativa} \\ \hline
Basal & \texttt{Random}, \texttt{Negligent} & Crítica (Alto sufrimiento / Alta mortalidad) \\
Heurística Simple & \texttt{Severity} & Moderada (Subóptima por miopía algorítmica) \\
Metaheurística & \texttt{Simulated Annealing}, \texttt{ILS} & Elevada (Convergencia a regiones óptimas) \\
Híbrida & \texttt{SA-ILS} & Superior (Máxima robustez y dominancia) \\ \hline
\end{tabular}
\end{table}

Es destacable que la variante híbrida (\texttt{SA-ILS}) capitalizó las fortalezas complementarias de ambos paradigmas, 
consolidando los mejores indicadores globales. 
Mientras que el esquema de enfriamiento de SA proporciona una flexibilidad matemática ideal para evadir 
óptimos locales en las fases tempranas de la simulación, el operador de Hill Climbing de 
ILS ejecuta una intensificación agresiva y determinista sobre las soluciones prometedoras.

\subsection{Dinámica Comparativa Inter-Políticas}

A pesar de la paridad estadística general entre ILS y SA, 
sus dinámicas de convergencia difieren según la naturaleza de la instancia. 
SA manifiesta una alta sensibilidad a la topología del espacio de soluciones, 
logrando configuraciones superiores en escenarios con transiciones de fase marcadas, 
a expensas de una mayor volatilidad temporal. Por su parte, ILS ofrece un comportamiento altamente estable y 
computacionalmente eficiente, idóneo para la inserción incremental de pacientes en tiempo real. 

La combinación de estas propiedades en la política híbrida mitiga la varianza residual de los métodos individuales. 
Esta propiedad de dominancia estocástica es crucial para su transferencia a entornos de triaje clínico real, 
donde la predictibilidad de la política ante fluctuaciones imprevistas de la demanda es tan prioritaria como la 
optimización del valor esperado del rendimiento.

\subsection{Análisis de Sensibilidad Paramétrica}

La robustez de los algoritmos se evaluó sometiendo al sistema a variaciones en las condiciones de contorno 
del entorno hospitalario (capacidad instalada y tasas de arribo). Los resultados demuestran que la ventaja 
competitiva de las metaheurísticas es inversamente proporcional a la holgura de los recursos del escenario:

\begin{equation}
\lim_{\text{Recursos} \to \infty} \Delta(\text{Rendimiento}) = 0
\end{equation}

Bajo escenarios de alta disponibilidad, el problema de secuenciación se trivializa y las trayectorias de todas 
las políticas convergen. 
No obstante, ante restricciones severas de acoplamiento médico —donde el sistema opera cerca del punto de saturación—, 
la capacidad de sintonización fina de ILS y la exploración estocástica de SA previenen el colapso del flujo hospitalario. 
Finalmente, se constató que la variante híbrida exhibe la menor elasticidad respecto al ajuste de hiperparámetros internos 
($max\_iterations$ y régimen de temperatura), 
consolidándose como la arquitectura más resiliente frente a la incertidumbre del modelo.

% =========================================================
% 9. Limitaciones y mejoras
% =========================================================
\section{Limitaciones y posibles mejoras}

Aunque los resultados obtenidos son prometedores, el sistema presenta varias limitaciones. 
En primer lugar, el comportamiento clínico de los pacientes se modela mediante reglas simplificadas y parámetros sintéticos, 
por lo que no representa completamente la complejidad de la evolución médica real. 
Asimismo, el dataset fue generado artificialmente con apoyo de un LLM, lo que puede introducir sesgos o 
inconsistencias respecto a datos clínicos reales.

Por otra parte, la evaluación de las soluciones depende de simulaciones estocásticas, 
por lo que una misma política puede producir resultados ligeramente distintos entre ejecuciones. 
Además, las metaheurísticas empleadas no garantizan alcanzar el óptimo global, sino aproximaciones de alta calidad.

Como líneas futuras de trabajo, se propone incorporar datos hospitalarios reales, modelar una mayor variedad de recursos y 
restricciones, incluir diferentes tipos de pacientes y niveles de atención, 
así como realizar análisis estadísticos sobre múltiples corridas para obtener intervalos de confianza. 
También sería interesante explorar técnicas más avanzadas de optimización, como algoritmos genéticos, 
búsqueda tabú, aprendizaje por refuerzo o enfoques híbridos que combinen simulación, 
optimización y modelos de lenguaje para apoyar la toma de decisiones clínicas.

% =========================================================
% 10. Conclusiones
% =========================================================
\section{Conclusiones}

En este trabajo se abordó el problema de priorización de casos clínicos bajo recursos limitados mediante 
una combinación de simulación basada en eventos discretos, técnicas de optimización y modelos de lenguaje. 
Se desarrolló un entorno de simulación capaz de representar la llegada de pacientes, la evolución de su estado de salud, 
el consumo de recursos hospitalarios y las decisiones de admisión bajo restricciones operacionales.

El uso de un LLM permitió enriquecer descripciones clínicas y generar automáticamente atributos relevantes para la simulación, 
facilitando la construcción de un conjunto de datos coherente con el dominio del problema. 
Sobre este escenario se evaluaron diversas políticas de priorización, desde estrategias simples hasta 
metaheurísticas de búsqueda.

Los resultados obtenidos muestran que las políticas basadas en optimización superan consistentemente a las heurísticas básicas, 
logrando salvar una mayor cantidad de pacientes y reducir el deterioro clínico acumulado. 
En particular, la combinación de \textit{Iterated Local Search} y \textit{Simulated Annealing} 
presentó el comportamiento más robusto entre las alternativas evaluadas.

Finalmente, el proyecto demuestra que la integración de inteligencia artificial y simulación constituye 
una herramienta viable para apoyar la toma de decisiones en entornos hospitalarios con recursos escasos, 
permitiendo analizar diferentes estrategias de asignación y estudiar su impacto antes de aplicarlas en la práctica.

\end{document}
